\partchapter{4}{第四部分\quad 总结}

\Needspace{85pt}

\section{可扩展之处}

以下方向为后续计划；当前最终演示已采用 PUF3，§2.3.4 的功能通过不等于下述可靠性与安全性目标已完成。

\begin{enumerate}
\def\labelenumi{\arabic{enumi}.}
\tightlist
\item
  长消息支持：当前 PT/AD 各 $\leq$16 B，经 AXI-Lite 字节 FIFO 传输；可改为共享缓冲或流式搬运，并补齐完整 1089 组板级 KAT 与多块吞吐测量。
\item
  器件绑定密钥：在 PUF3 已接入的基础上修正 RO 计数器时序，评估关闭派生输入回读通路、helper 持久化与失败处理，补充跨上电、温压、BER、熵及多器件统计后再评估密钥用途。
\item
  指令与算法：在现有 custom-0 状态指令上扩展 SP 800-232 中的 Ascon-Hash256/XOF128，并以 intrinsic 封装降低软件调用开销。
\item
  工程验证：补 PUF3 完整错误恢复矩阵、小时级稳定性、完整 ISA/异常矩阵与干净环境复现；已有 mcycle/PS 时基频率估测，仪器级测量另行开展。
\end{enumerate}

\Needspace{85pt}

\section{心得体会}

实现可运行与结论可复核是不同目标。标准 KAT、真实 CPU 仿真与实板测试各自覆盖不同问题，不能用某一层通过替代另一层验收；相同文件名也不能证明工件属于同一构建。因此，发布身份、逐文件哈希、source\_dirty 和源集选择应与原始样本一并保存。

性能比较首先是测量边界的比较。当前固定 16 B 输入下 MMIO 与 ISE 均降低了固件调用周期，但 PS 侧调用仍包含 Python 与通信开销，响应时间没有获得同倍数改善。保留这一区别，比仅展示最大的周期加速比更有解释力。

验证框架同样需要负向用例。旧 release 的错误长度、过期 epoch、重复请求及 FIFO overflow 门禁，和 PUF3 的自动锁定反例、未注册拒绝 HWKEY、坏 Tag 拒绝，分别提供有限但可复查的证据。它们不构成无限时稳定性、完整 ISA 合规或安全认证证明；版本变化后不能靠文字转移通过结论。
